% !TEX program = xelatex
% 그림(brief_summary.pdf)은 이 폴더에서 포함합니다.
\documentclass[10pt,a4paper]{article}

\usepackage{kotex}
\IfFontExistsTF{Malgun Gothic}{
  \setmainhangulfont{Malgun Gothic}
  \setsanshangulfont{Malgun Gothic}
}{
  \setmainhangulfont{Noto Serif CJK KR}
  \setsanshangulfont{Noto Sans CJK KR}
}
\usepackage{amsmath}
\usepackage{booktabs,tabularx,array}
\usepackage{graphicx}
\usepackage{xcolor}
\usepackage{geometry}
\usepackage[hidelinks]{hyperref}
\geometry{top=1.55cm,bottom=1.55cm,left=1.70cm,right=1.70cm}
\setlength{\parindent}{0pt}
\setlength{\parskip}{0.38em}
\setlength{\emergencystretch}{2em}
\renewcommand{\arraystretch}{1.15}
\renewcommand{\figurename}{그림}
\renewcommand{\tablename}{표}

\definecolor{navy}{HTML}{17324D}
\definecolor{teal}{HTML}{0F766E}
\definecolor{orange}{HTML}{C2410C}
\definecolor{pale}{HTML}{EEF6F5}
\definecolor{warn}{HTML}{FFF7ED}
\newcommand{\Status}[1]{\textsf{\footnotesize[#1]}}
\newcommand{\sensunit}{\mathrm{nV\,cm^{-1}\,Hz^{-1/2}}}
\newcommand{\summarybox}[1]{%
  \noindent\colorbox{pale}{%
    \parbox{\dimexpr\textwidth-2\fboxsep\relax}{#1}}}
\newcommand{\warnbox}[1]{%
  \noindent\colorbox{warn}{%
    \parbox{\dimexpr\textwidth-2\fboxsep\relax}{#1}}}

\graphicspath{{./}}
% 본 보고에 사용한 분석 산출값 (2026-07-24 재실행, 38점 온도 격자).
\providecommand{\SlopeOptimumTemperatureC}{43.00}
\providecommand{\SensitivityOptimumTemperatureC}{48.00}
\providecommand{\ModelToPPTScaleRoom}{2.06}
\providecommand{\GabesOptimumTransmissionPct}{3}
\providecommand{\PaperPSNOptimumC}{33}
\providecommand{\PaperOptimumTransmissionPct}{43}
\providecommand{\RbMeltingC}{39.30}
\providecommand{\AbsoluteSensitivityStatus}{PREDICTED}
\providecommand{\ExperimentalInputStatus}{PPT}

\begin{document}

{\color{navy}\LARGE\bfseries Rydberg 셀 가열에 따른 RF 감도 변화\par}
\vspace{2pt}
{\large 예비 시뮬레이션 및 주경민 선배님 실험 대조 보고}

\vspace{0.5em}
\begin{tabularx}{\textwidth}{>{\bfseries}p{2.7cm}X}
보고 대상 & 주경민 선배님, ``cell heating 이후 RF sensitivity 변화양상 확인'' (2026-07-15)\\
기준 논문 & G.\ Ju \textit{et al.}, ``Photon shot-noise-limited Rydberg-EIT electrometry,''
           arXiv:2606.04354v1 (선배님 제1저자, 부산대 문한섭 교수님 그룹)\\
시뮬레이션 & GABES --- 박수인이 개발 중인 원자광학 시뮬레이션 코드입니다.
             전체 코드는 \href{https://github.com/Shake2313/fwm-squeezing-app.git}{GitHub 저장소}에서 볼 수 있습니다.\\
작성 및 교정 & 원자광학 시뮬레이션(GABES) 실행 후 AI 지원으로 초안을 작성하고 세 차례 교정·교열했습니다.
             분석 2026년 7월 24일, 교정 2026년 8월 6일\\
증거 수준 & 발표자료 전사 \Status{PPT}, 선배님 논문 게재값 \Status{REFERENCE},
            시뮬레이션 계산 \Status{PREDICTED},
            미측정 입력 \Status{ASSUMED}/\Status{PENDING}\\
\end{tabularx}

\vspace{0.4em}
\summarybox{%
\textbf{\color{navy}보고 요약}\quad
셀을 가열해 Rb 밀도를 높였을 때 EIT 신호와 RF 최소 검출 전기장이 어떻게 변하는지 확인하여
후속 마이크로파 센싱의 운전 온도를 정하는 시험입니다. 현재 자료에서 읽을 수 있는 핵심은 세 가지입니다.
\textbf{(1)} PPT의 네 조건에서 측정 감도와 기재된 PSN 한계의 비는 0.95--1.12입니다. 두 값은
불확도 내에서 양립하지만, 이 비만으로 기술 잡음의 부재를 입증할 수는 없습니다. PPT에 기재된 PSN 제한
전기장 감도는 주석 온도 38\,°C에서 44\,°C 사이에 3.7배 증가합니다.
\textbf{(2)} GABES의 기울기 최적점 \SlopeOptimumTemperatureC\,°C와 PPT의 ``EIT 신호가 가장 뚜렷한''
조건에 병기된 42.5\,°C는 수치상 가깝습니다. 그러나 서로 다른 광학 지표를 비교하고 주석 온도도
미확정이므로, 이는 검증 결과가 아니라 후속 대조 후보입니다.
\textbf{(3)} 기준 논문은 PSN 한계가 \(\sim\PaperPSNOptimumC\,°C\)에서 최소가 된 뒤 고온에서
악화되는 거동을 제시합니다. 반면 GABES의 조건부 감도 최적점은 \SensitivityOptimumTemperatureC\,°C입니다.
이 불일치는 모델의 응답도와 검출 광자 수 사이 온도 의존 균형을 점검해야 함을 보여주지만,
투과 광량 항 하나를 원인으로 특정하지는 못합니다.}

\section*{1. 어떤 실험인지, 그리고 기준 논문과의 관계}

발표자료에 담긴 내용은 다음과 같습니다.

\begin{enumerate}
  \setlength{\itemsep}{0.15em}
  \item 세라믹 히터와 온도 센서로 증기 셀 가열 장치를 제작합니다. \Status{PPT}
  \item 상온 및 히터 설정값 40/45/50\,°C에서 EIT 신호를 비교합니다. 45\,°C 설정값에서
        신호가 가장 뚜렷하다고 기록되어 있습니다. \Status{PPT}
  \item 각 조건에서 1\,s 측정을 기준으로 최소 검출 전기장과 PSN 한계를 구합니다. PSN 주석에는
        설정값과 별개로 38/42.5/44\,°C가 ``??''와 함께 병기되어 있습니다. 이 값들의 물리적 의미와
        정확도는 확인되지 않았으므로 이하에서는 ``PPT 주석 온도''로 부릅니다. \Status{PPT}
\end{enumerate}

이 실험은 독립된 히터 검증이 아니라, 선배님 논문(arXiv:2606.04354v1)에서 확립한
\(\mathbf{{}^{85}\mathrm{Rb}}\) 계단형 EIT 원자 초헤테로다인 전기계의 \textbf{운전점 교정} 단계입니다.
그 논문은 상온(\(\sim20\,°C\)) 증기 셀에서 \(12.5(8)\,\sensunit\)의 감도(증기 셀 최고 기록,
PSN 한계 \(11.2\,\sensunit\)에 근접)를 보고했고, \textbf{셀 온도에 따른 PSN 곡선(Fig.\ 5(a))까지
이미 제시}하고 있습니다. 본 보고의 온도 대조는 그 곡선을 기준으로 삼습니다.
발표자 계획에도 ``07/16--24 온도에 따른 RF sensitivity 추가 확인''이 있어, 선배님 본인도 이 자료를
확정값이 아닌 1차 스캔으로 취급한 것으로 해석됩니다.

\begin{table}[h]
\centering
\small
\caption{PPT 전사값과 단순 계산한 비. 감도는 낮을수록 좋습니다. 주석 온도에는 ``??''가 붙어 있어
미확정입니다. 상온 감도의 중앙값은 논문과 같지만, 표의 불확도 \(\pm3.6\)은 PPT 전사값이며
논문값 \(12.5(8)\)의 불확도 \(\pm0.8\)과 다릅니다.}
\setlength{\tabcolsep}{4.4pt}
\begin{tabularx}{\textwidth}{@{}lcccc>{\raggedright\arraybackslash}X@{}}
\toprule
조건 & \shortstack{PPT 주석\\온도} & \shortstack{RF 감도\\\((\sensunit)\)} & PSN 한계 & 측정/PSN & 비고\\
\midrule
상온 & --- & \(12.5\pm3.6\) & 11.2(4) & 1.12 & PPT 전사; 논문값과 중앙값 동일\\
40\,°C 설정값 & 38\,°C?? & \(\mathbf{6.68\pm2.55}\) & 6.72(4) & 0.99 & 네 조건 중 중앙값 최저, \emph{측정 순서상 마지막}\\
45\,°C 설정값 & 42.5\,°C?? & \(17.6\pm2.2\) & 17.6(5) & 1.00 & EIT 신호가 가장 뚜렷함\\
50\,°C 설정값 & 44\,°C?? & \(23.8\pm6.5\) & 25.0(1) & 0.95 & 감도 악화\\
\bottomrule
\end{tabularx}
\end{table}

\section*{2. PPT 자체 정합성 점검 \Status{PPT}}

아래 비와 배수는 PPT 전사값에서 단순 계산했으며, 원시 데이터로 검증되지 않았으므로 증거 상태를
\Status{PPT}로 유지합니다.

\textbf{측정 감도와 기재된 PSN 한계는 불확도 내에서 양립합니다.} 네 조건의 측정/PSN 비는
0.95--1.12입니다. 동일한 광량·대역폭·보정 규약을 사용했다면 큰 초과 잡음의 증거는 보이지 않지만,
이 비만으로 기술 잡음의 부재를 입증하거나 검출기 RIN, 전자단 잡음, 환경 진동을 각각 배제할 수는 없습니다.

\textbf{PPT에 기재된 PSN 제한 전기장 감도도 고온 조건에서 악화됩니다.} 주석 온도 38\,°C와 44\,°C에
대응하는 값은 6.72에서 25.0\,\(\sensunit\)로 3.72배 증가합니다. 논문 Eq.\,(2)의 형태에 따르면
PSN 제한 감도는 대략 \(S_{\mathrm{PSN}}\propto(P_pT_p)^{-1/2}|\partial\alpha_0|^{-1}\)이므로,
검출 광자 수와 헤테로다인 응답 기울기를 함께 확인해야 합니다. 현재 전사값만으로 어느 항이 변했는지는
구분할 수 없습니다.

\warnbox{\small\textbf{재현을 위해 확인할 규약.}\quad 40\,°C와 50\,°C 조건에서는 중앙 측정값이
각각의 PSN 한계보다 약간 작습니다(0.99, 0.95). PSN 산정에 사용한 투과 광량, QE, RBW/VBW/ENBW,
평균 시간과 피크/RMS 규약을 원기록에서 확인해야 이 비를 정량 해석할 수 있습니다.}

\section*{3. 선배님 논문 온도 모델과의 대조 \Status{PREDICTED vs REFERENCE}}

이 대조가 본 보고의 핵심입니다. 선배님 논문 Fig.\ 5(a)는 셀 온도에 대한 PSN 한계 곡선을
제시하며, \textbf{최적점을 \(\sim\PaperPSNOptimumC\,°C\)}로, 그리고 \(\sim39\,°C\) 부근의 증기압
\textbf{불연속}(Rb 고체→액체 상전이)을 보고합니다. 현 시뮬레이션을 38점 온도 격자
(20--55\,°C)로 재실행해 PSN(T)와 투과율(T)을 산출했습니다.

\begin{tabularx}{\textwidth}{>{\bfseries\raggedright\arraybackslash}p{4.15cm}p{2.0cm}X}
\toprule
비교 항목 & 상태 & 값과 해석\\
\midrule
논문 PSN 최적 온도 & \Status{REFERENCE} & \(\sim\PaperPSNOptimumC\,°C\), Fig.\ 5(a)의 보고값\\
논문 증기압 분기 & \Status{REFERENCE} & \(\sim39\,°C\), Rb 고체--액체 상전이 부근\\
PPT 네 조건 중 감도 중앙값 최저 & \Status{PPT} & 40\,°C 설정값; 주석 온도 38\,°C??, 원시 데이터 미확인\\
GABES 조건부 감도 최적점 & \Status{PREDICTED} & \SensitivityOptimumTemperatureC\,°C\\
GABES 공진 투과율 & \Status{PREDICTED} & 48\,°C에서 \(\sim\GabesOptimumTransmissionPct\%\),
논문이 보고한 33\,°C에서 GABES 계산값 \(\sim\PaperOptimumTransmissionPct\%\)\\
\bottomrule
\end{tabularx}

논문 모델과 PPT 전사값은 모두 ``저온에서 개선된 뒤 고온에서 악화될 수 있다''는 정성적 그림과
양립합니다. 특히 42.5/44\,°C 주석 조건의 큰 PSN 값은 논문 곡선의 \textbf{고온 쪽 악화 구간
(감도 수치의 상승부)}과 같은 방향입니다. 다만 PPT에는 네 조건만 있고 주석 온도도 미확정이므로,
38\,°C를 실험 최적점으로 확정할 수는 없습니다.

\textbf{GABES와 기준 논문의 최적 온도는 어긋납니다.} GABES는 공진 투과율이
\(\sim\GabesOptimumTransmissionPct\%\)인 48\,°C에서도 응답도 이득이 검출 광자 수 감소를 상쇄한다고
계산합니다. 반면 논문이 보고한 최적 온도 33\,°C에서 GABES가 계산한 투과율은
\(\sim\PaperOptimumTransmissionPct\%\)입니다. 이 두 투과율은 모두 \Status{PREDICTED}이며 논문 게재값이
아닙니다. 따라서 불일치는 모델의 응답도와 광자 수 사이 온도 의존 균형이 지나치게 낙관적일 가능성을
보여주지만, 투과 광량 항 하나를 직접 원인으로 특정하지는 못합니다.

\warnbox{\small\textbf{온도와 상(phase)의 해석.}\quad 히터 설정값은 40/45/50\,°C이고 PPT 주석에는
38/42.5/44\,°C가 ``??''와 함께 병기되어 있습니다. 증기압 분기는 reservoir 또는 저온점 온도가
결정하므로 온도 로그 없이 각 조건의 고체·액체 분기를 확정할 수 없습니다. GABES는 설정값=유효 증기
온도=저온점 온도를 가정하고 CRC 고체/액체 증기압 분기를 적용했습니다.}

\section*{4. 시뮬레이션 광학/RF 응답 대조}

계산은 \({}^{85}\mathrm{Rb}\) \(5S_{1/2}\!\rightarrow5P_{3/2}\!\rightarrow40D_{5/2}\) EIT와
\(40D_{5/2}\!\rightarrow39F_{7/2}\) RF 결합을 사용했습니다(상태·기하·전력은 모두 선배님 논문 게재값).
각 온도에서 정상 상태 OBE 스펙트럼을 구하고, 40\,kHz finite-IF weak-SIG 선형 응답, RF 전이 쌍극자 모멘트,
탐침광 산탄 잡음을 전기장 ASD로 변환했습니다.

\begin{tabularx}{\textwidth}{>{\bfseries}p{3.4cm}X}
논문 게재 입력 & 탐침광 6\,$\mu$W, 결합광 30\,mW, 빔 지름 0.15\,mm, 셀 50\,mm,
  \(\Omega_{\mathrm{LO}}/2\pi=3.7\,\mathrm{MHz}\), \(\Omega_c/2\pi=3.0\,\mathrm{MHz}\),
  상태 \(40D_{5/2}\)/\(39F_{7/2}\),
  검출 양자효율 95\,\%, IF 40\,kHz \Status{REFERENCE}\\
조건부 가정 & 도플러 효과 미포함, 온도/밀도 의존 탈위상 0, 기술 잡음 0,
  RF 각도 인자 1, 히터 설정 온도=유효 증기 온도=저온점 온도, 검출 경로 효율 0.5 \Status{ASSUMED}\\
결과 상태 & 절대 감도 \Status{\AbsoluteSensitivityStatus}, 실험 입력 \Status{\ExperimentalInputStatus};
  실험값에 맞춘 피팅은 하지 않았습니다\\
\end{tabularx}

주요 광학 입력은 선배님 논문 게재값이지만, 절대 감도는 검출 경로 효율 0.5, RF 각도 인자 1,
기술 잡음 0과 온도 동일시 가정에도 의존합니다. 따라서 실험과 모델의 격차를 논문 입력 또는 모델 근사
한쪽에만 귀속할 수 없습니다.

\clearpage
\refstepcounter{figure}
\begin{center}
\IfFileExists{brief_summary.pdf}{%
  \includegraphics[width=0.68\textwidth,trim=0 0 406.63pt 22pt,clip]{brief_summary}\par\vspace{0.15em}
  \includegraphics[width=0.68\textwidth,trim=406.63pt 0 0 22pt,clip]{brief_summary}}{%
  \fbox{\parbox{0.9\textwidth}{\centering\vspace{1.0cm}
  그림 파일 \texttt{brief\_summary.pdf}을 찾을 수 없습니다.
  \vspace{1.0cm}}}}
\par\vspace{0.35em}
\parbox{0.95\textwidth}{\small\textbf{\figurename~\thefigure:}
GABES 온도 스윕과 PPT 전사값의 비교. 위는 서로 다른 두 모델 광학 지표이고,
아래의 빨간 점은 히터 설정값에 대한 PPT 전사값입니다. 회색 점선은 상온 중앙값에 맞춘 표시용
정규화이며 피팅 결과가 아닙니다.}
\end{center}

\textbf{광학 응답: 후속 대조 후보입니다.} 현 시뮬레이션의 스펙트럼 기울기는
\(\SlopeOptimumTemperatureC\,°C\)에서 최대입니다. PPT에서 ``EIT가 가장 뚜렷한'' 조건은 45\,°C
설정값이고, 여기에 42.5\,°C??가 병기되어 있습니다. 두 수치는 가깝지만 PPT의 ``뚜렷함''은 정성적
대비/진폭 관찰이고 모델 지표는 기울기이며, 모델의 EIT 대비와 기울기 최적 온도도 서로 다릅니다.
따라서 이를 광학 응답 경향의 검증으로 해석하지 않습니다.

\textbf{RF 감도 절댓값: 정량 일치가 확인되지 않았습니다.} 20\,°C에서 PPT 중앙값은 같은 온도의
GABES 값보다 약 \ModelToPPTScaleRoom 배 큽니다. 기존의 7.0배 비교는 38\,°C의 PPT PSN 값과
48\,°C의 GABES 최솟값을 나눈 것이어서 같은 온도·같은 물리량의 비교가 아니므로 교정본에서는
정량 근거로 사용하지 않습니다. 최적 온도 위치도 달라 단일 배율 보정의 타당성을 주장할 수 없습니다.

\begin{tabularx}{\textwidth}{>{\bfseries\raggedright\arraybackslash}p{4.0cm}X}
\toprule
현재 전사값에서 보이는 점 & 측정 감도와 기재 PSN 한계는 불확도 내에서 양립; 큰 초과 잡음의 증거는
  보이지 않지만 개별 기술 잡음은 배제하지 못함\\
\midrule
모델에서 검증할 항목 & 20\,°C GABES EIT 선폭 1.613\,MHz는 논문값 1.6\,MHz와 수치상 일치하나,
  잔류 도플러와 온도·밀도 의존 탈위상을 생략했으므로 고온 선폭·기울기·finite-IF 응답은 미검증;
  검출 경로 효율과 RF 각도 인자도 가정값\\
실험에서 확인할 항목 & 온도별 투과 전력과 IF 응답도; 저온점과 창면 응축; 히터 구조물의 37\,GHz 산란;
  검출기 RIN·전자단 잡음·대역폭 규약\\
공통 확인 항목 & 설정값, 유효 증기 온도와 저온점 온도의 관계; Zeeman·MW 정재파의 온도 상관 변화\\
\bottomrule
\end{tabularx}

\section*{5. 판단, 한계 및 다음 측정}

\summarybox{%
\textbf{\color{navy}판단}\quad
현재 자료로 말할 수 있는 것은 \textbf{``PPT 네 조건 중 40\,°C 설정값에서 감도 중앙값이 가장 낮았으므로
반복 측정한다''}까지입니다. 병기된 38\,°C??의 의미는 미확정이고, 이 조건은 측정 순서상 마지막이어서
시간 드리프트(잠금 품질, MW 전력, 정렬)와 온도 효과가 분리되지 않았습니다. 상온값은 논문 운전점과
중앙값이 같지만 불확도와 캠페인 대응이 불명확하므로 서로 다른 캠페인 사이의 유의도를 하나의
\(\sigma\) 값으로 제시하지 않습니다. 45\,°C 설정값은 EIT 신호가 가장 뚜렷했던 조건이지만 네 조건 중
RF 감도가 가장 낮은 조건은 아닙니다. GABES의 \(\SensitivityOptimumTemperatureC\,°C\)는 가정에 의존한
조건부 최솟값이므로 실험 목표 온도로 인용해서는 안 됩니다.}

\subsection*{권장 재측정 순서}

\begin{enumerate}
  \setlength{\itemsep}{0.22em}
  \item \textbf{측정 순서 무작위화·반복.} 상온을 이번 실험 구성에서 다시 측정해 기준점을 이 캠페인 안으로
        가져오고, 온도 로그를 기준으로 유효 증기 온도 33, 36, 38, 40, 42\,°C를 조밀하게 반복하며
        45, 48\,°C를 대조점으로 둡니다.
        논문 최적 \(\sim\PaperPSNOptimumC\,°C\) 부근을 포함하도록 저온 쪽을 넓히면 좋겠습니다.
        가열/냉각 양방향으로 드리프트와 온도 효과를 분리합니다.
  \item \textbf{광량과 응답도를 함께 기록.} 온도별 DC 투과 전력, 원시 EIT 주파수 스캔과
        IF 응답도(V per V/m)를 감도와 함께 남깁니다. 그래야 검출 광자 수 감소와 원자 응답 변화가
        PSN 및 총감도에 미치는 영향을 분리해 모델과 대조할 수 있습니다.
  \item \textbf{온도장 실측.} 히터 설정값 외에 양단 센서, 빔 부근의 유효 온도,
        reservoir와 저온점 후보 온도를 함께 기록합니다. PPT PSN 주석의 38/42.5/44\,°C가 무엇을 뜻하는지
        원기록에서 확정할 필요가 있습니다.
  \item \textbf{잡음 규약 고정.} 원시 EIT 스캔, RF 스윕, 검출기 암신호 트레이스와 PSD를 저장하고
        RBW, VBW, ENBW, 평균 시간, 피크/RMS 규약을 함께 기록합니다.
\end{enumerate}

\subsection*{보고 시 반드시 붙일 한계}

\begin{tabularx}{\textwidth}{>{\bfseries\raggedright\arraybackslash}p{3.35cm}p{2.05cm}X}
\toprule
항목 & 상태 & 의미\\
\midrule
PPT 수치 & \Status{PPT} & 원시 트레이스에서 재계산하지 않은 수동 전사값\\
상온 데이터 출처 & \shortstack{\Status{PPT}\\[-0.15em]\Status{REFERENCE}} & 중앙값 12.5와 PSN 11.2는 논문 운전점과 같지만,
  감도 불확도 \(\pm3.6\)은 PPT 전사값이고 논문값은 \(\pm0.8\); 캠페인 대응 미확인\\
측정 순서·드리프트 & \Status{PENDING} & 온도 효과와 시간 드리프트 미분리\\
시뮬레이션 스펙트럼 & \Status{PREDICTED} & 기준값에 고정된 정상 상태 OBE 결과\\
절대 RF 감도 & \Status{PREDICTED} & 기술 잡음 0인 조건부 PSN 기준; 실험 감도 아님\\
모델 입력 & \Status{REFERENCE} & 결합광 전력, 빔 지름, 셀 길이, LO/결합광 라비 주파수는
  선배님 논문 게재값\\
응답도·광자 수 균형 & \Status{PREDICTED} & 48\,°C 조건부 최적점의 원인을 한 항으로 분리하지 못함\\
온도장·저온점 & \Status{ASSUMED} & 설정값을 증기/저온점 온도와 동일시\\
Zeeman·RF 공간장 & \Status{PENDING} & 편광, 정재파, 3차원 장 미포함\\
원시 EIT/IF/PSD/온도 로그 & \Status{PENDING} & 정량 검증과 탈위상·잡음·온도 매핑 추정에 필요\\
\bottomrule
\end{tabularx}

\subsection*{최종 결론}

기준 논문의 온도 모델과 PPT 전사값은 고온에서 PSN 제한 감도가 악화될 수 있다는 정성적 그림과
양립합니다. 그러나 43\,°C 모델 기울기와 42.5\,°C?? PPT 관찰은 서로 다른 지표이므로 광학 검증으로
볼 수 없고, PPT 네 조건만으로 38\,°C를 실험 최적점으로 확정할 수도 없습니다. GABES의 조건부 감도
최적점 \(\SensitivityOptimumTemperatureC\,°C\)는 기준 논문의 \(\sim\PaperPSNOptimumC\,°C\)와 어긋나며,
이는 모델의 응답도·검출 광자 수·온도 매핑을 함께 점검해야 함을 뜻합니다. 다음 단계는 온도별 DC 투과,
원시 EIT와 IF 응답, PSD 및 저온점 로그를 같은 캠페인에서 확보해 실험 전사값과 모델을 병행 검증하는
것입니다. 현 단계의 적절한 표현은 ``검증 완료''가 아니라 \textbf{예비 시뮬레이션과 검증 과제 정리}입니다.

{\footnotesize
\textbf{재현 정보:} 본 보고의 정량값은 발표자료(2026-07-15) 전사 수치와 시뮬레이션 정상 상태 OBE 계산
결과를 구분했습니다. 기준 논문에서는 Fig.\ 5(a)의 보고된 랜드마크(최적
\(\sim\PaperPSNOptimumC\,°C\), 융점 \RbMeltingC\,°C)만 사용했으며 논문 곡선을 재구성하지 않았습니다.
표의 43\%와 3\% 투과율은 모두 GABES 38점 온도 격자 계산값입니다. 원본 발표자료와 인접 원본 그림은
수정하지 않았습니다. 분석 재실행일: 2026년 7월 24일; 3차 교정·교열일: 2026년 8월 6일.}

\end{document}
